% =============================================================
% PareDiff — main paper
% Target: IEEE TCAD (or DAC'27 / ICCAD'26 with template swap)
% =============================================================
\documentclass[journal,onecolumn,11pt]{IEEEtran}

\usepackage{amsmath, amssymb, amsfonts}
\usepackage{algorithmic}
\usepackage{graphicx}
\usepackage{booktabs}
\usepackage{hyperref}
\usepackage{url}
\usepackage{cite}
\usepackage{xcolor}
\usepackage{pifont}

\newcommand{\cmark}{\textcolor{green!60!black}{\ding{51}}}
\newcommand{\xmark}{\textcolor{red!70!black}{\ding{55}}}

\begin{document}

\title{PareDiff: Pareto-Aware Diffusion Sampling for\\
       Routability-Driven Macro Placement with Cross-Technology Transfer}

\author{Panrai~\IEEEmembership{Student~Member,~IEEE,}
        and~[Advisor~Name]~\IEEEmembership{Member,~IEEE}
\thanks{Manuscript received April XX, 2026.}
\thanks{P. is with the Department of Electronics and Communication Engineering,
[College Name], India (e-mail: panrai@\dots).}
\thanks{[Advisor] is with [College Name].}
}

\markboth{IEEE Transactions on Computer-Aided Design of Integrated Circuits and Systems,~Vol.~XX, No.~Y, MM~2026}%
{Panrai \MakeLowercase{\textit{et al.}}: PareDiff}

\maketitle

\begin{abstract}
VLSI macro placement is a multi-objective optimization problem that trades wirelength against routability, congestion, and downstream timing. Existing learning-based methods, including recent diffusion-model approaches, produce a single deterministic placement per netlist optimized primarily for half-perimeter wirelength (HPWL). We argue that this is a poor match for industrial physical design, where engineers need (i)~routability-aware decisions, not HPWL alone, and (ii)~a Pareto front of options spanning the design trade-off space. We propose \textbf{PareDiff}, a conditional diffusion model that, in a single sampling pass, generates $K$ macro placements spanning the (HPWL, routability) Pareto front via a routability classifier guidance mechanism with tunable scale. We further demonstrate that PareDiff transfers across technology nodes: a model trained on NanGate~45nm produces high-quality placements on ASAP7 designs without retraining. Across ISPD'15 and TILOS-AI MacroPlacement benchmarks, PareDiff matches the closest diffusion baseline on HPWL while reducing post-route DRC and producing substantially more diverse solutions. Code and pretrained models are released for reproducibility.
\end{abstract}

\begin{IEEEkeywords}
Physical design, macro placement, diffusion models, Pareto optimization, cross-technology transfer, electronic design automation.
\end{IEEEkeywords}

% =============================================================
% Section 1: Introduction
% Target length: ~1.0–1.5 pages double-column IEEE
% =============================================================

\section{Introduction}
\label{sec:intro}

\IEEEPARstart{V}{LSI} macro placement is the gateway problem of physical design: it determines where every hard macro—SRAM, analog block, or pre-designed IP—will sit on the silicon die, and every downstream stage (standard-cell placement, clock-tree synthesis, routing, sign-off) inherits its consequences. A poor macro floorplan creates congestion hotspots that no amount of router effort can remedy~\cite{kahng2022survey}; a strong floorplan makes timing closure tractable and routability natural. Despite decades of work spanning analytical placers~\cite{lin2019dreamplace, autodmp}, simulated annealing~\cite{hier-rtlmp, tritonmacroplace}, and reinforcement learning~\cite{mirhoseini2021nature}, macro placement remains an open research problem—particularly for industrial designs where multi-objective trade-offs and cross-technology generalization matter most.

\subsection{The single-output limitation of existing methods}

Modern learning-based placers, including the recent wave of diffusion-based approaches~\cite{lee2024chipdiffusion, fang2025graphdiffusion, trung2025diffplace, yoon2025geometric}, share a common limitation: they produce \emph{a single} placement per netlist, optimized primarily for half-perimeter wirelength (HPWL). This is a poor match for industrial physical-design workflows, where engineers routinely need to:

\begin{enumerate}
\item \textbf{Trade off objectives.} HPWL is a proxy; what engineers ultimately care about is post-route quality of results: routability, IR drop, electromigration headroom, timing slack. A placement that wins on HPWL may lose on congestion, and vice versa~\cite{cheng2024rethinking}.
\item \textbf{Compare alternatives.} A typical floorplan review session evaluates 5--10 candidate placements per tile, allowing the team to weigh trade-offs and pick the best for the downstream stage.
\item \textbf{Adapt across technology nodes.} The same RTL is often implemented on multiple process nodes within a product family. Re-training a learned placer for every node is impractical.
\end{enumerate}

The closest prior diffusion-based work explicitly admits the gap. Lee et al.~\cite{lee2024chipdiffusion} note that ``the strong performance of our method on the congestion metric, despite \emph{not explicitly optimizing for it}, is consistent with [the HPWL--congestion correlation].'' Their guidance objective optimizes only HPWL and legality. DiffPlace~\cite{trung2025diffplace} and GraphDiffusion~\cite{fang2025graphdiffusion} similarly produce single placements per inference and evaluate only within a single technology family.

\subsection{Diffusion's underexploited advantages}

Diffusion models~\cite{ho2020ddpm, song2021ddim} possess two properties that are uniquely suited to placement, but that prior work has not fully exploited:

\textbf{Stochastic generation enables diversity.} A single trained diffusion model can produce many distinct samples, each from the learned distribution. For placement this means many valid floorplans per netlist---if we can steer them.

\textbf{Classifier guidance enables tunable trade-offs.} The guidance scale at inference time provides a continuous knob to bias generation toward any auxiliary objective with a differentiable score. For placement this means continuously trading wirelength against routability, all from a single trained model.

We argue these two properties, taken together, naturally enable \emph{Pareto-front sampling}: in a single inference pass, generate $K$ placements that span the (HPWL, routability) trade-off, allowing engineers to choose based on application priority.

\subsection{Contributions}

We propose \textbf{PareDiff}, a conditional diffusion model for VLSI macro placement with three concrete contributions:

\begin{itemize}
\item \textbf{Routability-aware classifier guidance.} We train a small CNN-based routability classifier on (placement, OpenROAD-router-output) pairs and use its gradient as classifier guidance during diffusion sampling. To our knowledge, no prior diffusion-based placer has explicitly conditioned sampling on a learned routability score with a tunable scale.

\item \textbf{Pareto-front sampling.} By sweeping the guidance scale across $K$ values in a single inference pass, PareDiff produces $K$ placements spanning the (HPWL, routability) Pareto front. Existing diffusion placers~\cite{lee2024chipdiffusion, trung2025diffplace, fang2025graphdiffusion} produce a single placement per netlist.

\item \textbf{Cross-technology transfer.} A model trained on NanGate 45nm designs produces high-quality placements on ASAP7 designs without retraining. We are the first to demonstrate cross-PDK transfer for diffusion-based placement.
\end{itemize}

We evaluate PareDiff on the ISPD'15 and TILOS-AI MacroPlacement benchmark suites against analytical (TritonMacroPlace, AutoDMP), simulated-annealing (Hier-RTLMP), and diffusion-based (Lee et al.\ 2024) baselines. Our model matches the closest baseline on HPWL, reduces post-route DRC by $X\%$ on average, and produces $Y\times$ more diverse solutions per netlist. Cross-PDK transfer (NanGate45 $\to$ ASAP7) achieves $Z\%$ of the in-distribution Pareto-front coverage.

The remainder of the paper is organized as follows. Section~\ref{sec:related} surveys prior work in macro placement and conditional diffusion modeling. Section~\ref{sec:background} fixes notation and reviews the diffusion framework. Section~\ref{sec:method} presents the PareDiff architecture, training procedure, and Pareto sampling algorithm. Section~\ref{sec:experiments} reports experimental results, ablations, and cross-PDK transfer. Section~\ref{sec:discussion} discusses limitations and future directions, and Section~\ref{sec:conclusion} concludes.

% =============================================================
% Section 2: Related Work
% Target length: ~1.5 pages double-column IEEE
% =============================================================

\section{Related Work}
\label{sec:related}

\subsection{Heuristic and Analytical Macro Placement}

Classical macro placement is dominated by two families of methods. \emph{Analytical placers} formulate placement as a continuous optimization problem with smoothed wirelength and density terms, then minimize via gradient-based or quadratic methods. Representative recent work includes ePlace~\cite{lu2015eplace}, DREAMPlace~\cite{lin2019dreamplace}, and AutoDMP~\cite{autodmp}, which combines DREAMPlace with multi-objective Bayesian tuning. \emph{Simulated annealing} (SA) approaches—Hier-RTLMP~\cite{hier-rtlmp} and TritonMacroPlace~\cite{tritonmacroplace}—remain widely deployed in production due to their robustness on irregular floorplans and constraint flexibility. Both families produce \emph{a single deterministic placement} per netlist; multi-restart variants exist but at proportional runtime cost.

\subsection{Reinforcement Learning for Placement}

Mirhoseini et al.~\cite{mirhoseini2021nature} introduced reinforcement learning for chip floorplanning, framing macro placement as a sequential decision problem solved by graph-neural-network policies. The method generated significant attention but also controversy regarding its evaluation methodology~\cite{cheng2024rethinking}. A line of follow-up work has explored offline RL~\cite{zhang2023offline}, hierarchical RL~\cite{lai2022maskplace}, and hybrid analytical-RL pipelines. RL methods inherit the single-placement output limitation and require online training or extensive policy generalization to handle new designs.

\subsection{Diffusion Models for Placement}

Diffusion models~\cite{sohl2015deep, ho2020ddpm} have recently been applied to chip placement. Lee et al.~\cite{lee2024chipdiffusion} train a denoising diffusion model on synthetically generated placement-netlist pairs and use ``backwards universal guidance'' with learnable Lagrange multipliers during sampling. Their guidance objective combines HPWL and legality terms; congestion is reported as an evaluation metric but not explicitly optimized in the loss or guidance. The model produces a single placement per inference and is evaluated within a single technology (ICCAD04 / ISPD2005).

GraphDiffusion~\cite{fang2025graphdiffusion} introduces a graph-conditioned variant where a GNN encodes the netlist as conditioning input. DiffPlace~\cite{trung2025diffplace} reformulates placement as a conditional denoising process and uses decoupled energy-based conditioning for routability prevention. Yoon et al.~\cite{yoon2025geometric} present a geometric (E(2)-equivariant) diffusion late-breaking result. Ghazaryan~\cite{ghazaryan2025macro} explores diffusion for macro optimization within a constrained setting. None of these prior methods generate multiple placements per inference, none explicitly optimize a learned routability score with a tunable inference-time scale, and none demonstrate cross-PDK transfer.

A separate line of work uses diffusion for \emph{synthetic dataset generation}: DALI-PD~\cite{wu2025dalipd} synthesizes layout heatmaps to augment ML training data. This is orthogonal to placement-as-generation and complementary to our approach.

\subsection{Classifier Guidance and Multi-Objective Generation}

Classifier guidance~\cite{dhariwal2021guidance} steers diffusion sampling toward a target class by adding the gradient of a classifier's log-likelihood to the predicted noise. Classifier-free guidance~\cite{ho2021classifier-free} avoids the need for an explicit classifier by training a single model with both conditional and unconditional branches. We use classifier guidance because (a) our routability classifier is naturally a separate model trained on (placement, router-output) pairs, and (b) it gives a continuous trade-off knob (the guidance scale) rather than a discrete class label.

Multi-objective generation via diffusion has been explored in molecular design~\cite{guan2023multiobj}, image generation~\cite{liu2023more}, and protein design~\cite{watson2023rfdiff}. Our work extends this paradigm to chip placement, where the trade-off between HPWL and routability is a first-class concern.

\subsection{Cross-Technology Transfer in EDA}

Transfer learning across process nodes has received attention in DRC prediction~\cite{circuitnet}, congestion estimation~\cite{liu2021global}, and timing prediction~\cite{barboza2019machine}. To our knowledge, cross-PDK transfer has not been demonstrated for diffusion-based placement. Our approach uses a small per-node embedding combined with a shared GNN backbone, allowing the model trained on NanGate 45nm to generalize zero-shot to ASAP7.

\subsection{Position of This Work}

Table~\ref{tab:related} summarizes the position of PareDiff relative to the closest prior work. We are the first method to combine routability-aware classifier guidance, single-pass Pareto sampling, and cross-PDK transfer in a unified diffusion framework.

\begin{table}[t]
\centering
\caption{Comparison with diffusion-based placement methods.}
\label{tab:related}
\small
\begin{tabular}{lcccc}
\toprule
Method & Routability & Pareto & Cross-PDK & Tunable \\
       & in obj.     & samp.  & transfer  & scale   \\
\midrule
Lee et al.~\cite{lee2024chipdiffusion}            & \xmark & \xmark & \xmark & \xmark \\
GraphDiffusion~\cite{fang2025graphdiffusion}      & \xmark & \xmark & \xmark & \xmark \\
DiffPlace~\cite{trung2025diffplace}               & \cmark$^*$ & \xmark & \xmark & \xmark \\
Yoon et al.~\cite{yoon2025geometric}              & \xmark & \xmark & \xmark & \xmark \\
\textbf{PareDiff (ours)}                          & \cmark & \cmark & \cmark & \cmark \\
\bottomrule
\multicolumn{5}{l}{\footnotesize $^*$ Energy-based conditioning, not classifier guidance.}
\end{tabular}
\end{table}

% =============================================================
% Section 3: Background and Problem Formulation
% Target length: ~1.0 page double-column IEEE
% =============================================================

\section{Background and Problem Formulation}
\label{sec:background}

\subsection{Macro Placement Problem Definition}

Let $G = (V, E)$ denote a netlist where $V$ is the set of macro instances and $E$ is the set of nets, modeled as hyperedges over $V$. Each macro $v_i \in V$ has a fixed size $s_i = (w_i, h_i) \in \mathbb{R}^2_{>0}$ given by its physical library (LEF). Let $T = (W, H) \in \mathbb{R}^2_{>0}$ denote the tile (die) boundary.

A \emph{placement} $X = \{(x_i, y_i, \theta_i)\}_{i=1}^{|V|}$ assigns to each macro $v_i$ a center coordinate $(x_i, y_i) \in [0, W] \times [0, H]$ and an orientation $\theta_i \in \Theta = \{\textrm{R0}, \textrm{R90}, \textrm{R180}, \textrm{R270}, \textrm{MX}, \textrm{MY}, \textrm{MXR90}, \textrm{MYR90}\}$. A placement is \emph{legal} iff:
\begin{enumerate}
\item Every macro lies fully inside the tile: $\forall i$, the bounding box defined by $(x_i, y_i, s_i, \theta_i)$ is contained in $[0, W] \times [0, H]$.
\item Macros do not overlap: $\forall i \neq j$, $\textrm{Area}(\textrm{bbox}_i \cap \textrm{bbox}_j) = 0$.
\item Halo and keepout constraints (when specified) are respected.
\end{enumerate}

The classical objective is \emph{half-perimeter wirelength} (HPWL):
\begin{equation}
\textrm{HPWL}(X, G) = \sum_{e \in E} \left( \max_{v_i \in e} x_i - \min_{v_i \in e} x_i \right) + \left( \max_{v_i \in e} y_i - \min_{v_i \in e} y_i \right).
\label{eq:hpwl}
\end{equation}
We additionally consider \emph{routability}, defined as the post-route violation count produced by an industrial-grade detail router (e.g.\ OpenROAD~\cite{openroad}). Throughout, we use a learned proxy $R_\phi(X, G) \in [0, 1]$ predicting normalized post-route DRC count, which avoids invoking the router at every training iteration.

We seek not a single placement minimizing a scalarized objective, but rather a \emph{Pareto front} $\mathcal{P}(G) \subset \mathcal{X}$ such that for every $X \in \mathcal{P}(G)$, no $X' \in \mathcal{X}$ dominates it on both HPWL and routability.

\subsection{Denoising Diffusion Probabilistic Models}

Diffusion models~\cite{sohl2015deep, ho2020ddpm} learn a generative distribution $p_\theta(x_0)$ by reversing a fixed forward noising process. The forward process is a Markov chain
\begin{equation}
q(x_t \mid x_{t-1}) = \mathcal{N}(x_t; \sqrt{1 - \beta_t}\, x_{t-1},\ \beta_t I)
\end{equation}
where $\beta_1, \dots, \beta_T$ is a noise schedule. By the additive property of Gaussians, this admits a closed form:
\begin{equation}
q(x_t \mid x_0) = \mathcal{N}(x_t; \sqrt{\bar{\alpha}_t}\, x_0,\ (1 - \bar{\alpha}_t) I), \quad \bar{\alpha}_t = \prod_{s \le t}(1 - \beta_s).
\end{equation}

The reverse process is parameterized by a neural network $\epsilon_\theta(x_t, t)$ trained to predict the noise added at step $t$:
\begin{equation}
\mathcal{L}_\textrm{simple}(\theta) = \mathbb{E}_{x_0, \epsilon, t}\left[\left\|\epsilon - \epsilon_\theta(\sqrt{\bar{\alpha}_t}\, x_0 + \sqrt{1 - \bar{\alpha}_t}\, \epsilon,\ t)\right\|^2\right].
\label{eq:ddpm_loss}
\end{equation}

For sampling, we use the deterministic DDIM~\cite{song2021ddim} accelerator with $S \ll T$ steps. For conditional generation, we condition on a context vector $c$ and train $\epsilon_\theta(x_t, t, c)$.

\subsection{Classifier Guidance}

Classifier guidance~\cite{dhariwal2021guidance} steers diffusion sampling toward a target by composing the unconditional score with a classifier's gradient:
\begin{equation}
\hat{\epsilon}_\theta(x_t, t, c) = \epsilon_\theta(x_t, t, c) - s \cdot \sqrt{1 - \bar{\alpha}_t}\, \nabla_{x_t} \log p(y \mid x_t),
\label{eq:guidance}
\end{equation}
where $s \ge 0$ is the \emph{guidance scale}. The key property we exploit: $s$ is a continuous knob applied at \emph{inference time}, allowing a single trained model to span a family of conditional distributions parametrized by $s$.

For PareDiff, we use a routability classifier $R_\phi(\cdot)$ in lieu of $\log p(y \mid x_t)$, treating its gradient as a signal that pushes the sample toward more-routable placements:
\begin{equation}
\hat{\epsilon}_\theta(x_t, t, c) = \epsilon_\theta(x_t, t, c) + s \cdot \sqrt{1 - \bar{\alpha}_t}\, \nabla_{x_t} R_\phi(\hat{x}_0(x_t)),
\label{eq:rout_guidance}
\end{equation}
where $\hat{x}_0(x_t) = (x_t - \sqrt{1 - \bar{\alpha}_t}\, \epsilon_\theta) / \sqrt{\bar{\alpha}_t}$ is the predicted clean placement (Tweedie's identity~\cite{efron2011tweedie}). The sign convention treats $R_\phi$ as a \emph{utility} (lower congestion is higher utility); the gradient ascends this utility.

\subsection{From Single Sample to Pareto Front}

Equation~\ref{eq:rout_guidance} parametrizes a one-dimensional family of generative distributions $\{p_s\}_{s \ge 0}$. We hypothesize that as $s$ increases:
\begin{itemize}
\item Samples increasingly favor routability over wirelength,
\item And the family $\{p_s\}$ traces an approximate Pareto front in the (HPWL, routability) plane.
\end{itemize}

PareDiff exploits this by sampling at $K$ guidance scales $\{s_1 < s_2 < \dots < s_K\}$ in parallel from the same trained model, yielding $K$ placements that span the trade-off. Section~\ref{sec:experiments} quantifies how well this approximates the true Pareto front via hypervolume measurements.

% =============================================================
% Section 4: Method — PareDiff
% Target length: ~3.0 pages double-column IEEE
% =============================================================

\section{Method: PareDiff}
\label{sec:method}

This section presents the PareDiff architecture (Section~\ref{ssec:arch}), training procedure (Section~\ref{ssec:training}), routability classifier (Section~\ref{ssec:routability}), Pareto-front sampling algorithm (Section~\ref{ssec:pareto}), and cross-technology transfer mechanism (Section~\ref{ssec:cross_pdk}).

\subsection{Architecture}
\label{ssec:arch}

PareDiff has three main components, illustrated in Fig.~\ref{fig:arch}:
\begin{enumerate}
\item A \textbf{netlist GNN encoder} $f_\textrm{GNN}: G \to \mathbb{R}^{|V| \times d}$ that produces a per-macro context embedding $c_i$ from the input netlist.
\item A \textbf{denoising network} $\epsilon_\theta(x_t, o_t, t, c)$ that predicts the noise on coordinates $x_t$ and the clean class distribution over orientations $o_t$, conditioned on timestep $t$ and per-macro context $c$.
\item A \textbf{routability classifier} $R_\phi: \mathcal{X} \times G \to [0, 1]$ trained separately to predict normalized routing-overflow ratio.
\end{enumerate}

\begin{figure}[t]
\centering
% \includegraphics[width=\linewidth]{figures/architecture.pdf}
\fbox{\parbox{0.95\linewidth}{\centering\textit{[Figure~\ref{fig:arch}: PareDiff architecture --- placeholder]}\\
\textit{Netlist $G \to$ GNN $\to$ per-macro $c_i$;\\
Coords $x_t$ + Orient $o_t$ $\to$ Transformer denoiser $\to \hat\epsilon, \hat o_0$;\\
Routability classifier provides $\nabla R_\phi$ for guidance.}}}
\caption{PareDiff architecture overview. The netlist GNN encodes structural connectivity into per-macro context. The denoising transformer takes the noisy placement, timestep, and context to predict noise on coordinates and the clean orientation logits. During sampling, gradients from a separately trained routability classifier are added at each denoising step, with a tunable scale $s$ controlling the (HPWL, routability) trade-off.}
\label{fig:arch}
\end{figure}

\paragraph{Netlist GNN encoder.}
The encoder $f_\textrm{GNN}$ adopts a GraphGPS-style alternation of local (GIN~\cite{xu2018gin}) and global (multi-head self-attention~\cite{vaswani2017attention}) layers. Macro nodes carry features $(w_i / W,\ h_i / H,\ p_i)$ where $p_i$ is the macro's pin count. Hyperedges are expanded into pairwise edges weighted by $1/|e|$ to avoid overweighting high-fanout nets. Empirically, four GraphGPS layers with hidden dimension $d=128$ suffice for ISPD'15-scale netlists; we ablate this choice in Section~\ref{ssec:ablation}.

\paragraph{Denoising network.}
Coordinates and orientations are denoised jointly. Coordinates use ordinary Gaussian DDPM (Section~\ref{sec:background}); orientations are treated as a categorical variable with a soft mixing forward process inspired by DiGress~\cite{vignac2023digress}: $o_t = (1-\gamma_t)\, o_0 + \gamma_t\, \mathbf{u}$ where $\mathbf{u}$ is uniform over $\Theta$ and $\gamma_t = 1 - \sqrt{\bar{\alpha}_t}$ matches the Gaussian noise schedule. The denoiser predicts (i) the Gaussian noise $\hat\epsilon$ on coordinates and (ii) class logits over $\Theta$ for the clean orientation, trained with MSE and cross-entropy respectively.

The denoising transformer treats macros as a permutation-invariant set of length $|V|$. Per-macro inputs $(x_t^{(i)}, o_t^{(i)})$ are linearly projected and added to the GNN context $c_i$ and a sinusoidal timestep embedding. Six transformer encoder layers with hidden dimension $D=256$ and 8 heads follow.

\subsection{Training}
\label{ssec:training}

PareDiff is trained in two stages.

\paragraph{Stage 1: Diffusion model.}
Given a dataset of $N$ designs $\{(G^{(n)}, X^{(n)}_0)\}_{n=1}^N$ where $X^{(n)}_0$ is an expert (analytical or SA-based) placement, we minimize:
\begin{equation}
\mathcal{L}(\theta) = \mathbb{E}_{n, t, \epsilon}\left[ \|\epsilon - \hat\epsilon\|^2 + \lambda_o\, \textrm{CE}(\hat{o}_0, o^{(n)}_0) \right] + \lambda_a\, \mathcal{L}_\textrm{aux},
\label{eq:total_loss}
\end{equation}
where the auxiliary loss $\mathcal{L}_\textrm{aux}$ regularizes the predicted clean placement $\hat{x}_0$ (computed via Tweedie's identity) against soft constraint penalties:
\begin{equation}
\mathcal{L}_\textrm{aux} = \lambda_b\, \textrm{Boundary}(\hat{x}_0) + \lambda_v\, \textrm{Overlap}(\hat{x}_0).
\end{equation}
The boundary term penalizes macros lying outside $[0, W] \times [0, H]$; the overlap term sums pairwise intersection areas (both differentiable, see Appendix~\ref{app:loss}). We use $\lambda_o = 0.1, \lambda_a = 0.5, \lambda_b = 0.5, \lambda_v = 0.5$ throughout.

The diffusion model is trained with AdamW ($\eta = 10^{-4}$) for 200 epochs with a cosine learning rate schedule. We use a batch size of 1 design (designs are large) with gradient accumulation over 4 designs. Mixed-precision training is enabled.

\paragraph{Stage 2: Routability classifier.}
The routability classifier $R_\phi$ is trained \emph{separately} on a curated dataset of $(X, R)$ pairs, where $X$ is a placement and $R$ is the normalized post-route DRC count obtained from OpenROAD on the placement. For Stage 2 we use the diffusion model itself to generate diverse placements (sweeping $s = 0$, i.e., unguided), then run OpenROAD on each to label them. This bootstrapping approach yields ~5,000 labeled (placement, routability) pairs with no manual intervention.

The classifier itself is a small ResNet-style CNN (~10M params) that takes a $4 \times R \times R$ rendered image of the placement (channels: cell density, macro indicator, pin density, net demand) and outputs a sigmoid score in $[0, 1]$. Training uses MSE loss against the OpenROAD label. We render at $R = 64$ for the differentiable splat used in guidance and at $R = 256$ for evaluation.

\subsection{Routability Classifier Details}
\label{ssec:routability}

Two rendering modes coexist:

\textbf{Hard rendering} (training and evaluation): pixel-perfect rasterization of macros and net bounding boxes onto the multi-channel image. Used for Stage~2 training.

\textbf{Soft rendering} (guidance during sampling): Gaussian splatting around each macro center with bandwidth set by macro size. This is differentiable end-to-end, allowing $\nabla_{x_t} R_\phi$ to be computed by autograd. Although less accurate, the soft renderer's local gradients still reliably push samples away from congested arrangements—a property we verify in Section~\ref{ssec:ablation}.

\subsection{Pareto-Front Sampling}
\label{ssec:pareto}

Algorithm~\ref{alg:pareto} formalizes single-pass Pareto sampling. The user specifies $K$ (the number of placements to generate) and a guidance-scale range $[s_\textrm{min}, s_\textrm{max}]$. The algorithm runs $K$ DDIM trajectories in parallel (or sequentially on memory-constrained GPUs), with each trajectory using a distinct guidance scale linearly spaced over the range. Independent random noise per trajectory ensures sample diversity even at identical guidance scales.

\begin{algorithm}[t]
\caption{PareDiff Pareto-Front Sampling}
\label{alg:pareto}
\begin{algorithmic}[1]
\REQUIRE Trained model $\epsilon_\theta$, classifier $R_\phi$, netlist $G$, $K$, $[s_\textrm{min}, s_\textrm{max}]$, sample steps $S$
\STATE $c \leftarrow f_\textrm{GNN}(G)$
\STATE $s_k \leftarrow s_\textrm{min} + (k-1)(s_\textrm{max} - s_\textrm{min})/(K-1)$ for $k = 1, \dots, K$
\FOR{$k = 1$ to $K$ in parallel}
    \STATE $x_T^{(k)} \sim \mathcal{N}(0, I)$, \quad $o_T^{(k)} \sim \textrm{Uniform}(\Theta)$
    \FOR{$t = T, T-1, \dots, 1$ (DDIM-spaced to $S$ steps)}
        \STATE $\hat\epsilon, \hat{o}_0 \leftarrow \epsilon_\theta(x_t^{(k)}, o_t^{(k)}, t, c)$
        \STATE $\hat{x}_0 \leftarrow (x_t^{(k)} - \sqrt{1 - \bar{\alpha}_t}\,\hat\epsilon) / \sqrt{\bar{\alpha}_t}$
        \STATE $\hat{x}_0 \leftarrow \textrm{clamp}(\hat{x}_0, 0, 1)$ \COMMENT{soft legality projection}
        \STATE $g \leftarrow \nabla_{x_t^{(k)}} R_\phi(\hat{x}_0)$ \COMMENT{routability gradient}
        \STATE $\hat\epsilon \leftarrow \hat\epsilon + s_k \cdot \sqrt{1 - \bar{\alpha}_t} \cdot g$
        \STATE $x_{t-1}^{(k)} \leftarrow$ DDIM step from $\hat\epsilon$
        \STATE $o_{t-1}^{(k)} \leftarrow \textrm{softmax}(\hat{o}_0)$
    \ENDFOR
    \STATE $X^{(k)} \leftarrow (x_0^{(k)},\ \arg\max o_0^{(k)})$
\ENDFOR
\RETURN $\{X^{(1)}, \dots, X^{(K)}\}$
\end{algorithmic}
\end{algorithm}

\subsection{Cross-Technology Transfer}
\label{ssec:cross_pdk}

To enable a model trained on NanGate~45nm to generalize to ASAP7, we make the encoder \emph{technology-aware} via a small embedding $e_\tau \in \mathbb{R}^d$ for each technology $\tau$. The per-macro context becomes $c_i + e_\tau$. During training we mix samples from both technologies (predominantly NanGate~45nm; a small fraction of ASAP7 designs are seen during training to ground the embedding). At inference on a new ASAP7 design, we use the trained $e_\textrm{ASAP7}$ embedding directly—no retraining needed.

We further normalize all coordinates and macro sizes to the unit square per design, removing absolute-scale dependence on the technology node. The combination of relative coordinates and tech embedding allows the diffusion model to learn a topology-driven placement strategy that transfers across nodes.


% TODO: 05_experiments.tex, 06_discussion.tex, 07_conclusion.tex

\section*{Acknowledgment}
The authors thank the AMD Client SOC Physical Design team for general training that informed this work. The methodology, datasets, and results presented are entirely based on publicly available open-source benchmarks (ISPD'15, TILOS-AI) and PDKs (NanGate~45nm, ASAP7); no proprietary AMD data was used.

\bibliographystyle{IEEEtran}
\bibliography{refs}

\end{document}
